\chapter{Literature Review}

\section{Importance of mint as a medicinal and spice plant}

\subsection{Mint (\textit{Mentha spp.})}
Mint species are widely cultivated as medicinal and aromatic plants due to their economic value and diverse applications in the food, pharmaceutical, and cosmetic industries \textcite{Zheljazkov2020}. Among these, \textit{Mentha piperita}, \textit{M. spicata}, and \textit{M. crispata} are particularly important because of their high biomass production and essential oil content \textcite{Mahendran2021}.

The medicinal importance of mint has been recognized since ancient times. Historical records indicate that Nicander of Colophon used mint to create a repellant stench to ward off snakes \textcite{Silva2020}. This highlights the long-standing role of mint in traditional medicine and its relevance in modern therapeutic applications.

Mint is also valued for its essential oil composition. Peppermint leaves contain approximately 1.2--3.9\% essential oils, comprising over 300 identified compounds. The terpenic class is the most dominant, including about 52\% monoterpenes and 9\% sesquiterpenes, while other components such as aldehydes (9\%), alcohols (6\%), miscellaneous compounds (8\%), lactones (7\%), and aromatic hydrocarbons (9\%) are present in smaller proportions \textcite{Silva2020}.

Botanically, mint (\textit{Mentha spp.}) is an aromatic, predominantly perennial herb belonging to the family Lamiaceae. It is widely recognized for its dual importance as both a medicinal and culinary plant. Members of the genus \textit{Mentha} have been used since ancient civilizations and continue to play a significant role in modern medicine, food industries, and traditional health systems worldwide \textcite{Silva2020}.

\subsection{Medicinal Importance of Mint}
Mint has a long and well-established history of medicinal use, especially for the treatment of digestive disorders. Previous studies indicate that \textit{Mentha} species possess carminative, antispasmodic, antiemetic, antimicrobial, antioxidant, anti-inflammatory, and insecticidal properties. These medicinal properties are largely due to essential oils and phenolic compounds such as menthol, menthone, flavonoids, and phenolic acids.

Peppermint (\textit{Mentha piperita}) is one of the most extensively studied species and is commonly used for gastrointestinal ailments, including indigestion, flatulence, colic, nausea, and irritable bowel syndrome. Mint essential oils have also demonstrated antimicrobial activity against various bacteria and fungi, supporting their traditional use in treating infections and preserving food. Additionally, mint oils and menthol are widely used in pharmaceutical preparations, oral hygiene products, cosmetics, and aromatherapy due to their soothing, cooling, and antiseptic effects.

Historical evidence further highlights the medicinal relevance of mint, showing its use across ancient Babylonian, Greek, Roman, Islamic, and medieval European medical practices. According to \textcite{Silva2020}, ``Mints have been used since ancient Babylon, but it was in Classical Antiquity that their medical uses flourished.''

\subsection{Importance of Mint as a Spice Plant}
In addition to its medicinal value, mint is an important spice plant widely used in culinary practices across different cultures. The leaves, which can be consumed fresh or dried, the primary edible parts of the plant ``Either fresh or dried leaves of mint can be used for culinary purposes. Most often fresh mint is chosen over dried mint.''

Fresh leaves are preferred for their stronger aroma and flavor. Mint leaves are characterized by a warm, sweet, and refreshing taste with a distinctive cooling aftereffect, making them a popular flavoring agent in a wide range of foods and beverages, including sauces, jellies, confectionery, syrups, ice creams, and traditional dishes. Mint plays an important role in Middle Eastern, Indian, European, and African cuisines, where it is valued for both flavor enhancement and digestive benefits. Mint essential oils are also widely used in the food industry as natural flavoring agents and preservatives due to their antimicrobial properties.

\subsection{Economic and Cultural Significance}
The combined medicinal and spice value of mint has made it an economically important 
crop in many countries worldwide. The cultivation and processing of mint, 
both fresh and dried leaves, is vital to agricultural and industrial sectors. 
Its continued use in medicine, food, cosmetics, and perfumery underscores the 
enduring importance of mint as a versatile plant of global significance. 
According to \textcite{Silva2020}, ``They are still presently used for numerous 
purposes, including non-medicinal, which makes them economically relevant herbs.''

\section{Botanical and economic importance of \textit{Mentha piperita}, 
\textit{Mentha spicata}, and \textit{Mentha crispata}}

\begin{itemize}
    \item \textit{Mentha piperita} (peppermint) is one of the most economically important mint species due to its high menthol and menthone. According to \textcite{Schmidt2009}, ``menthol (40.7\%) and menthone (23.4\%) in essential oil.'' Menthol provides a cooling sensation and is widely used in pharmaceutical preparations, oral hygiene products, cosmetics, and food flavoring.

    \item \textit{Mentha spicata} (spearmint) contains a high amount of carvone in its essential oil. According to \textcite{Pino1998},``Identifies carvone as the major compound (65\%) in spearmint oil.'' Because of its mild and sweet flavor, spearmint is extensively used in the perfumery industries, flavoring industry, particularly in chewing gum, toothpaste, beverages, and confectionery products.

    \item \textit{Mentha crispata} (curly mint) is considered a variety of spearmint and is distinguished by its curled or wrinkled leaves. It is commonly used in culinary applications such as salads, herbal teas, and garnishes, and also contributes to essential oil production.

    \item \textit{Mentha suaveolens} (apple mint) has essential oils that are rich in compounds like piperitenone oxide, which make it biologically active \textcite{Kokkini2013}. The oil has bioactive monoterpenes that kill germs and have other health benefits, so it can be used in food and medicine \textcite{Hussain2010}. Mint species, such as \textit{Mentha suaveolens}, are also important for agriculture because they produce a lot of biomass and essential oils \textcite{Zheljazkov2013}.
\end{itemize}        

\section{Use of peat moss in indoor and greenhouse mint production}
Peat moss is widely used in greenhouse mint production due to its favorable physical properties; however, its extraction is associated with significant environmental impacts, including carbon emissions and peatland degradation \textcite{Cleary2021}. Peatlands store approximately one-third of the world’s soil carbon, but when they are drained for horticultural use, they shift from carbon sinks to major sources of CO\textsubscript{2} emissions. Degraded peatlands contribute nearly 5\% of global CO\textsubscript{2} emissions, exceeding the combined emissions from international aviation and shipping. 

In addition, peat extraction leads to the destruction of important habitats and threatens biodiversity in sensitive ecosystems such as Canada’s Hudson Bay Lowlands. Sustainability is also a major concern, as peat forms very slowly (approximately 1 mm per year), while extraction removes soil at a much faster rate, making it effectively a non-renewable resource. These environmental concerns have prompted regulatory actions, including the United Kingdom’s plan to phase out peat use in horticulture by 2030 and European Union initiatives aimed at reducing peat dependency.

Consequently, biochar has been proposed as a sustainable alternative growing media component due to its porous structure, chemical stability, and potential to improve substrate properties \textcite{Dumroese2018}. Therefore, despite its horticultural advantages, the environmental impacts of peat moss highlight the urgent need for alternative growing media.

\subsection{Overview and Economic Importance}
According to \textcite{yaghobvand2022evaluation}, species of the genus \textit{Mentha} are of high economic value due to their active and aromatic substances and are used as raw materials in food, cosmetics, health, beverage, and pharmaceutical industries. The increasing demand for high-quality essential oils requires a better understanding of the factors influencing plant growth and physiological traits.

\section{Growth Requirements and Physiology of Mint}

\subsection{Mentha piperita (Peppermint)}

\subsection*{Botanical Characteristics}
\textit{Mentha × piperita} is a rhizomatous, upright perennial herbaceous plant. According to the \textcite{rhs2026mentha}, peppermint is a hybrid, originally treated as a species but now known to be a cross between \textit{Mentha aquatica} (watermint) and \textit{Mentha spicata} (spearmint). It produces upright stems with dark green leaves and reddish stems that have a characteristic warm, spicy scent.

\subsection*{Climate and Environmental Requirements}
Research by \textcite{virzo1980adaptability} demonstrated that height growth of \textit{Mentha piperita} is better at 44\% than 100\% daylight. They further reported that at 14\% daylight, growth is limited, probably by lack of photosynthate, and no flowering occurs. Their transfer experiments between light and shade indicated that peppermint responds very effectively to irradiance also in a late stage of the life cycle with changes of specific leaf area and chlorophyll content.

Regarding temperature, environmental stresses including extreme temperatures can diminish yield and increase susceptibility to diseases. In some climates, as noted by the \textcite{umd2026growing}, summer heat may require partial shade; in subtropical regions, summer heat can reduce plant vigor unless some shade is provided.

According to the \textcite{mobot2026piperita}, the plant requires adequate moisture and adapts well to medium to wet soil conditions.

\subsection*{Physiological Responses to Nutrient Stress}
\textcite{srivastava1994relationship} investigated the relationship between photosynthetic carbon metabolism and essential oil biogenesis in peppermint under manganese stress. They found that under manganese deficiency, $CO_{2}$ exchange rate, total chlorophyll, total assimilatory area, plant dry weight, and essential oil yield were significantly reduced, whereas chlorophyll a/b ratio, leaf area ratio, and leaf stem ratio significantly increased. In leaves of manganese-deficient plants, $CO_{2}$ incorporation into the primary metabolic pool and essential oil were significantly lower.

In related research, \textcite{srivastava1993relation} examined peppermint under iron deficiency. They reported that $CO_{2}$ exchange rate, stomatal conductance, oil content, chlorophyll content, and chlorophyll a/b ratio were significantly affected. Growth parameters such as leaf area ratio, specific leaf weight, and leaf stem ratio were also influenced by iron stress.

\subsection*{Soil and Nutrient Requirements}
According to the \textcite{umd2026growing}, peppermint grows best in rich, moist soils but adapts to a wide range of soil conditions except dry ones. Soil quality, particularly the presence of organic matter and essential nutrients, is crucial for robust growth.

\subsection{Mentha spicata (Spearmint)}

\subsection*{Botanical Characteristics}
According to the \textcite{rhs2026spicata}, spearmint (\textit{Mentha spicata}) produces pointed, slightly crinkled leaves that are lighter green in color than peppermint, with the whole plant having a characteristic sweet smell. It is an aromatic perennial spreading by creeping rhizomes.

\subsection*{Physiological Responses to Water Stress}
\textcite{marino2019evaluation} conducted a field experiment on spearmint under different irrigation regimes in Southern Italy. They found that increasing levels of water stress affected later developing leaves, plant water status, and net photosynthesis from the beginning of stress. Photosynthesis limitation played a critical role during the harvest period. Under severe water stress, the crop restricted water losses by modulating stomatal closure and showed lowered mesophyll conductance. They further reported that irrigation treatments did not affect the concentration of organic compounds, while the yield of essential oils was negatively affected by water stress due to reduced crop growth in terms of total and leaf biomass, leaf area index, and crop height.

\subsection*{In Vitro Culture Responses}
Research by \textcite{tisserat2004influence} on physical parameters affecting spearmint cultures in vitro showed that the type of physical support greatly influenced culture growth and morphogenesis. They reported that mint shootlets grown on liquid medium only produced four-fold more fresh weight than shootlets grown on agar medium. Positive correlations occurred between culture vessel capacity and growth, morphogenesis, and carvone levels. An automated plant culture system allowed for the production of approximately fifteen-fold greater biomass compared to culture tubes.

Further research by \textcite{tisserat2008growth} confirmed these findings, demonstrating that carvone was only produced from cultured spearmint shoots and was absent in either roots or callus cultures. They also found that thenumber of media culture immersions greatly affected growth, morphogenesis, and secondary metabolism, with twelve immersions per day being optimal for growth and morphogenesis.

\subsection*{Comparative Growth Studies}
\textcite{yaghobvand2022evaluation} evaluated five mint species in an aeroponic system under greenhouse conditions. Their results showed that \textit{M. spicata} had the highest leaf-to-stem ratio. Most yield-attributes traits including leaf area, leaf fresh and dry weight, root fresh weight, shoot fresh weight, and total plant dry weight, as well as relative water content and photosynthesis rate, were higher in the second harvest than the first.

\subsection{Mentha crispata (Curly Mint)}

\subsection*{Botanical Characteristics and Taxonomy}
According to the \textcite{rhs2026crispa}, \textit{Mentha spicata} var. \textit{crispa} (curly spearmint) is an aromatic perennial reaching 60 cm in height, spreading by creeping rhizomes. The leaves are bright green, broad, and crinkled. Whorls of tiny, pale purple flowers appear on terminal spikes in summer.

The \textcite{rhs2026crispa} lists several synonyms for this plant, including:
\begin{itemize}
    \item \textit{Mentha spicata} 'Crispa'
    \item \textit{Mentha viridis} var. \textit{crispa}
    \item \textit{Mentha crispa} L. (1753)
\end{itemize}

\subsection*{Cultivation Requirements}
According to the \textcite{rhs2026crispa}, curly mint should be grown in any moist soil including chalk, clay, loam, or sand. Soil should be moist but well-drained, though plants tolerate poorly-drained conditions. The plant requires moist conditions but well-drained soil and should be protected from excess winter wet. It performs best in full sun to partial shade and is suitable for west, south, or east-facing positions.

Regarding growth management, the \textcite{rhs2026crispa} notes that this plant may have the potential to become a nuisance. They recommend restricting root run by planting in deep containers and plunging into soil, or growing in small, contained beds. Plants grown in pots benefit from dividing every few years.

\subsection{Mentha suaveolens (Apple Mint)}

\subsection*{Botanical Characteristics}
\textit{Mentha suaveolens} Ehrh., commonly known as apple mint, is a perennial herbaceous species belonging to the Lamiaceae family. It is characterized by soft, pubescent leaves and a distinctive fruity aroma. The species spreads through rhizomes and is well adapted to temperate environments, typically occurring in moist habitats \textcite{Magdy2025}.

\subsection*{Climate and Environmental Requirements}
Species within the genus \textit{Mentha}, including \textit{Mentha suaveolens}, generally thrive in temperate climates with moderate temperatures and sufficient water availability. Environmental factors such as salinity and extreme climatic conditions can significantly influence plant growth and productivity. Under optimal conditions, adequate moisture and moderate irradiance support vegetative development and essential oil production. However, stress conditions such as high salinity can reduce biomass accumulation and limit growth performance \textcite{Ahmadi2023, Magdy2025}.

\subsection*{Physiological Responses to Water and Salinity Stress}
Salinity stress has a pronounced effect on the physiological processes of \textit{Mentha suaveolens}. Increasing concentrations of NaCl significantly reduce fresh and dry biomass while altering mineral nutrient balance. In particular, sodium (Na$^+$) accumulation increases, whereas potassium (K$^+$) uptake decreases, leading to ionic imbalance and reduced metabolic efficiency \textcite{Kasrati2014}. Under moderate stress, plants may exhibit adaptive responses such as increased essential oil content; however, severe stress results in decreased growth and physiological activity \textcite{Ahmadi2023}.

\subsection*{Soil and Nutrient Requirements}
\textit{Mentha suaveolens} performs best in moist, fertile soils with adequate organic matter. Soil nutrient availability plays a crucial role in plant growth and essential oil synthesis. Adverse soil conditions, particularly high salinity, negatively affect nutrient uptake by disrupting ion homeostasis, especially through increased Na$^+$ levels and reduced K$^+$ availability. Therefore, maintaining balanced soil fertility and proper moisture conditions is essential for optimal cultivation \textcite{Kasrati2014, Magdy2025}.

\subsection*{Propagation}
For propagation, the \textcite{rhs2026crispa} recommends division in spring or autumn.

\section*{Comparative Physiology Across Mint Species}

\subsection*{Aeroponic System Performance}
\textcite{yaghobvand2022evaluation} conducted a comprehensive evaluation of five \textit{Mentha} species in an aeroponic system, revealing species-specific physiological responses.

Top-performing species: \textit{M. aquatica} showed superior performance compared to other species in yield-related attributes, including leaf number, leaf area, and fresh and dry weight of plants, shoots, and leaves. Among the remaining species, the highest shoot dry weight and total plant dry weight were observed in \textit{M. piperita}. Although \textit{Mentha suaveolens} was not the top-performing species, it demonstrated moderate biomass production, indicating good adaptability to controlled aeroponic conditions.

Physiological measurements: \textit{M. piperita} exhibited the highest fresh root weight, carotenoid and chlorophyll content, photosynthetic rate, and CO$_2$ exchange rate at the stomatal level. In comparison, \textit{Mentha suaveolens} showed moderate physiological activity, consistent with findings that this species maintains stable photosynthetic performance under varying environmental conditions, though generally lower than highly productive species such as \textit{M. piperita} \textcite{Kasrati2014, Ahmadi2023}.

Leaf-to-stem ratio: \textit{M. spicata} had one of the highest leaf-to-stem ratios, indicating a favorable allocation of biomass toward photosynthetically active tissues. \textit{Mentha suaveolens} also exhibited a relatively balanced leaf-to-stem ratio, contributing to efficient light capture and moderate biomass accumulation.

Harvest timing: Most yield-related traits, including leaf area, fresh and dry leaf weight, root fresh weight, shoot fresh weight, total plant dry weight, relative water content, and photosynthetic rate, were higher in the second harvest compared to the first. This trend was also observed in \textit{Mentha suaveolens}, suggesting that the species benefits from regrowth cycles and improved physiological efficiency over successive harvests.

Overall, while \textit{Mentha suaveolens} does not consistently outperform species such as \textit{M. aquatica} or \textit{M. piperita} in aeroponic systems, it demonstrates stable growth, moderate physiological performance, and good adaptability, making it suitable for controlled environment cultivation.

\section{Abiotic stress in mint cultivation with emphasis on salt stress}
Salt stress is a major abiotic constraint affecting plant growth and productivity, particularly in controlled cultivation systems where salinity can accumulate in the root zone, as reported by \textcite{Hasanuzzaman2021}. In medicinal and aromatic plants, salt stress negatively affects growth, photosynthesis, and secondary metabolite production, as described by \textcite{AcostaMotos2017}.

\section{Growth requirements and physiology of mint species}
Mint species require well-aerated substrates, adequate water supply, and balanced nutrient availability to achieve optimal growth and biomass production \textcite{Zheljazkov2020}. Physiological differences among mint species influence their tolerance to environmental stresses \textcite{Singh2022}.

\section{Salt stress effects on medicinal and spice plants}
Salt stress induces osmotic stress and ion toxicity, leading to reduced growth and yield in medicinal and spice plants \textcite{Shahid2020}. These effects are often more pronounced in container-grown crops due to limited root volume \textcite{Hasanuzzaman2021}.

\section{Physiological responses of plants to salt stress}

\subsection{Chlorophyll content and photosynthesis}
Salinity stress commonly leads to a reduction in chlorophyll content and photosynthetic efficiency as a result of chloroplast damage and stomatal limitation \textcite{Kalaji2021}. These effects directly contribute to reduced biomass accumulation under saline conditions \textcite{MunnsTester2008}.

\subsection{Water relations and osmotic adjustment}
Plants exposed to salt stress experience impaired water uptake and reduced leaf water potential, which can trigger osmotic adjustment mechanisms such as the accumulation of compatible solutes \textcite{AcostaMotos2017, Hussain2020}.

\section{Biochar effects on substrate properties and salinity stress}
Biochar application has been shown to improve substrate water-holding capacity, cation exchange capacity, and nutrient retention, which may alleviate the negative effects of salinity stress on plants \textcite{Akhtar2021}. Additionally, biochar can reduce sodium availability in the root zone, thereby mitigating ion toxicity under saline conditions \textcite{Joseph2021}.

\section{Research gap and justification}
The preceding sections establish that mint is an economically important crop, peat moss has significant environmental limitations, biochar offers promise as a sustainable alternative, and salt stress poses a real threat to greenhouse production. However, a critical knowledge gap remains: we do not know whether biochar can help mint plants tolerate salt stress. Specifically, there is limited research on the combined effects of biochar amendment and salt stress on different \textit{Mentha} species. No studies have systematically compared multiple mint species under these conditions to determine whether responses are species-specific or whether biochar's benefits vary across genotypes. 

Understanding these interactions is essential for developing evidence-based recommendations for growers seeking sustainable solutions to salt stress in mint production. This gap provides the rationale for the present study.

\section{Aim and objectives of the study}
The main aim of this study was to evaluate the effects of biochar as a peat alternative on the salt stress responses of three mint species. The specific objectives were to:

\begin{itemize}
    \item Assess growth responses of different mint species under salt stress
    \item Evaluate the influence of biochar-amended growing media on salt stress tolerance
    \item Compare physiological responses among mint species,
    \item Identify interactions between growing media composition and salinity stress.
\end{itemize}